\begin{trapbox}
\textbf{误区一：公式不显含 $a$，就认为任意固定 $b$ 和角度都能实现。}

已存在的焦点三角形，确实可用 $b,\theta$ 直接算面积；
但椭圆上的可实现范围还受到 $c$ 的限制：
\[
 \boxed{S\le cb,\qquad\tan(\theta/2)\le c/b}.
\]
\begin{center}\begin{tikzpicture}[font=\small,>=stealth]
\draw[black!30,rounded corners=2pt] (0,0) rectangle (14.4,8.2);
\node[anchor=west,font=\small\bfseries] at (.5,7.72) {同样 $b=3$，并不意味着同样的顶角都能实现};
\begin{scope}[shift={(3.6,4.2)},scale=.45]
\draw[blue,very thick] (0,0) ellipse (4 and 3);
\draw[orange!85!black,very thick] (-2.6457513111,0)--(0,3)--(2.6457513111,0)--cycle;
\fill (-2.6457513111,0) circle (3pt);\fill (2.6457513111,0) circle (3pt);
\end{scope}
\node at (3.6,1.6) {$a=4,\ b=3$};
\node at (3.6,2.25) {最高面积 $3\sqrt7<12$};
\begin{scope}[shift={(10.8,4.2)},scale=.45]
\draw[blue,very thick] (0,0) ellipse (5 and 3);
\draw[orange!85!black,very thick] (-4.0000000000,0)--(0,3)--(4.0000000000,0)--cycle;
\fill (-4.0000000000,0) circle (3pt);\fill (4.0000000000,0) circle (3pt);
\end{scope}
\node at (10.8,1.6) {$a=5,\ b=3$};
\node at (10.8,2.25) {最高面积 $12$};

\node[anchor=west] at (.7,.55) {同尺度；若要求 $\tan(\theta/2)=4/3$，左图不可能，右图恰在短轴端点实现};
\end{tikzpicture}\end{center}

两椭圆都有 $b=3$。若指定 $\tan(\theta/2)=4/3$，面积形式上给出 $12$；
左图最大面积只有 $3\sqrt7<12$，故该角不存在，右图才恰好实现。
这说明\textbf{“公式不含 $a$”不等于“任何 $a$ 都允许相同的 $\theta$”}。

\medskip
\textbf{误区二：把顶角误看成焦点处的角，或误写 $\tan\theta$。}

$\theta=\angle F_1PF_2$，顶点在 $P$；应写
$S=b^2\tan(\theta/2)$，\textbf{不是} $b^2\tan\theta$。
图中即使 $\theta$ 为钝角，$\theta/2$ 仍处于 $(0,90^\circ)$，半角正切为正。

\medskip
\textbf{误区三：把面积系数错记成 $a^2$ 或 $c^2$。}

由 $(a-c)(a+c)=a^2-c^2=\boxed{b^2}$，所以正确系数是\textbf{短半轴的平方}。
快速自检：取 $P=(0,b)$，底乘高 $S=cb$，且 $\tan(\theta/2)=c/b$，
只有 $b^2\tan(\theta/2)$ 给出 $cb$。

\medskip
\textbf{误区四：长轴端点依然是一个非退化三角形。}

当 $P=(\pm a,0)$ 时三点共线，$S=0$、$\theta=0$；
公式数值可以连续延伸，但不能对这样的退化图形使用内切圆或非退化夹角计算。

\medskip
\textbf{误区五：只反解半角，不验面积和角度是否取得到。}

反解后还要验证 $S\le cb$。例5中形式计算得 $16$，但 $16>12=cb$，故点不存在。

\medskip
\textbf{误区六：觉得顶角不能为钝角，或者认为所有椭圆都能出现钝角。}

看前面的三图对照：$c<b$ 时全为锐角；$c=b$ 时最大为直角；$c>b$ 时才存在钝角。

\end{trapbox}
\clearpage
\begin{trapbox}
\textbf{误区七：旋转、平移后继续生搬 $S=c|y|$。}
\begin{center}\begin{tikzpicture}[font=\small,>=stealth]
\draw[black!30,rounded corners=2pt] (0,0) rectangle (14.4,9.1);
\node[anchor=west,font=\small\bfseries] at (.5,8.6) {椭圆转过 $30^\circ$：高要垂直于焦点连线，不能直接读 $|y|$};
\begin{scope}[shift={(5.4,3.7)},scale=.7]
\draw[->,black!35] (-5.3,0)--(5.7,0) node[right] {$x$};
\draw[->,black!35] (0,-3.7)--(0,5.2) node[above] {$y$};
\begin{scope}[rotate=30]
\draw[blue,very thick] (0,0) ellipse (5 and 3);
\draw[black!60,dashed] (-5.3,0)--(5.3,0);
\fill[orange!12] (-4,0)--(0,3)--(4,0)--cycle;
\draw[orange!85!black,very thick] (-4,0)--(0,3)--(4,0)--cycle;
\draw[violet,very thick] (0,0)--(0,3);
\draw[violet] (0,.3)--(.3,.3)--(.3,0);
\fill (-4,0) circle (3pt);\fill (4,0) circle (3pt);\fill (0,3) circle (3pt);
\end{scope}
\node[right=6pt] at ({-2*sqrt(3)},-2) {$F_1$};
\node[below=6pt] at ({2*sqrt(3)},2) {$F_2$};
\node[above left=8pt] at (-1.5,{3*sqrt(3)/2}) {$P$};
\node[below right=5pt] at (0,0) {$O$};
\node[violet,rotate=-60] at (-1.2,1) {$d=3$};
\end{scope}
\node[anchor=west,align=left] at (10,5.4) {$a=5,b=3,c=4$\\[12pt]$\displaystyle |y_P|=\frac{3\sqrt3}{2}$\\[12pt]$d=3\ne |y_P|$\\[12pt]$S=cd=12$};
\node at (7.2,.6) {橙色底边仍长 $2c=8$；紫色垂线才是高。旋转不改变真实面积。};
\end{tikzpicture}\end{center}

图中纵坐标只量到水平 $x$ 轴，紫色高却垂直于倾斜的焦点连线；两者不同。
面积必须用 $S=cd$。只有焦点连线恰为 $x$ 轴时，才有 $d=|y|$；
中心为 $(h,k)$ 且长轴水平时，$d=|y-k|$。
半角式 $S=b^2\tan(\theta/2)$ 始终是几何关系，平移、旋转后仍成立。

\medskip
\textbf{误区八：拿椭圆结论直接套用双曲线或其他三角形。}

本公式用到了\textbf{椭圆第一定义 $PF_1+PF_2=2a$}，并结合 $b^2=a^2-c^2$。
双曲线第一定义是距离之差的绝对值等于 $2a$，不能直接替换使用本结论。
\end{trapbox}
\clearpage
